\subsection{Applications}

The Iterator pattern is extensively utilized across modern software engineering and systems programming to provide unified, decoupled traversal across heterogeneous data structures and streaming resources~\cite{gof1994, refactoring_guru_iterator}:

\begin{itemize}
    \item \textbf{Hierarchical DOM Traversal:} Modern web rendering engines use iterators to traverse recursive DOM trees without exposing internal node linkage.
    \item \textbf{Infinite Scrolling and Pagination:} Mobile and web feeds leverage iterators to paginate remote API data on demand, keeping RAM usage low.
    \item \textbf{Offline Sync Queues:} Mobile systems iterate sequentially through persistent task logs to replay operations after reconnecting to a network.
\end{itemize}

\subsubsection{Relationship with Other Design Patterns}

\paragraph{Iterator vs. Visitor}
\begin{table}[htbp]
\centering
\caption{Architectural Comparison: Iterator vs. Visitor}
\label{tab:iterator_vs_visitor}
\begin{tabular}{@{}p{3.5cm} p{5.5cm} p{5.5cm}@{}}
\toprule
\textbf{Dimension} & \textbf{Iterator Pattern}~\cite{gof1994} & \textbf{Visitor Pattern}~\cite{gof1994} \\ \midrule
\textbf{Core Objective} & Access and sequence traversal without exposing container internals. & Adding new operations to class hierarchies without modifying their source code. \\ \midrule
\textbf{Target Elements} & Primarily homogeneous or structured collections. & Heterogeneous object structures and composite hierarchies. \\ \midrule
\textbf{Dispatch Model} & Single dispatch (\texttt{it->next()}). & Double dispatch (\texttt{element->accept(visitor)}). \\ \midrule
\textbf{Traversal Control} & Explicitly driven by the Iterator cursor and client consumption loop. & Unconstrained by the pattern; traversal can be delegated to the Object Structure, an Iterator, or the Visitor itself. \\ \bottomrule
\end{tabular}
\end{table}

\noindent \textbf{Synergy:} An Iterator can traverse a complex composite hierarchy sequentially, passing each node to a Visitor that executes polymorphic operations without altering the tree structure~\cite{gof1994}.

\paragraph{Iterator vs. Chain of Responsibility}

\begin{table}[htbp]
\centering
\caption{Architectural Comparison: Iterator vs. Chain of Responsibility}
\label{tab:iterator_vs_cor}
\begin{tabular}{@{}p{3.5cm} p{5.5cm} p{5.5cm}@{}}
\toprule
\textbf{Dimension} & \textbf{Iterator Pattern}~\cite{gof1994} & \textbf{Chain of Responsibility Pattern}~\cite{gof1994} \\ \midrule
\textbf{Core Objective} & Traverses elements sequentially across a collection without exposing its internal structure. & Passes a request along a dynamic chain of potential handlers until one processes it. \\ \midrule
\textbf{Flow Control} & Sequential and deterministic; visits every element in order unless explicitly terminated by the client. & Conditional and delegated; execution passes down the handler link until a condition or termination criteria is met. \\ \midrule
\textbf{Decoupling Target} & Decouples the consumer algorithm from the underlying data structure layout. & Decouples the sender of a request from its potential receivers/handlers. \\ \midrule
\textbf{State Management} & Tracks iteration state (cursors, indices, tree pointers) explicitly. & Handlers maintain references only to their immediate successor in the chain. \\ \bottomrule
\end{tabular}
\end{table}

\noindent \textbf{Synergy (Processing Pipelines):}
The two patterns frequently collaborate in data processing pipelines, middleware stacks, and stream filtration~\cite{refactoring_guru_iterator}:
\begin{itemize}
    \item The \textbf{Iterator} acts as the \textit{data source pump}, fetching items one-by-one from complex, memory-efficient data structures without exposing internal node layouts.
    \item The \textbf{Chain of Responsibility} acts as the \textit{processing pipeline}, routing each item produced by the iterator through a series of discrete, loosely coupled validation, filtering, or transformation stages.
\end{itemize}

\begin{lstlisting}[language=C++, caption={Synergy between Iterator and Chain of Responsibility}]
#include <iostream>
#include <memory>
#include <string>
#include <vector>

struct Track {
    std::string title;
    int durationSec;
    bool isExplicit;
};

// --- CHAIN OF RESPONSIBILITY: Handler Interface ---
class TrackFilter {
protected:
    std::shared_ptr<TrackFilter> nextHandler;

public:
    virtual ~TrackFilter() = default;

    void setNext(std::shared_ptr<TrackFilter> next) {
        nextHandler = next;
    }

    virtual bool process(const Track& track) {
        if (nextHandler) {
            return nextHandler->process(track);
        }
        return true; // Passed all pipeline stages
    }
};

// Concrete Handler 1: Filter explicit content
class ExplicitFilter : public TrackFilter {
public:
    bool process(const Track& track) override {
        if (track.isExplicit) {
            std::cout << "[Blocked - Explicit]: " << track.title << "\n";
            return false;
        }
        return TrackFilter::process(track);
    }
};

// Concrete Handler 2: Filter out short tracks (< 30 seconds)
class DurationFilter : public TrackFilter {
public:
    bool process(const Track& track) override {
        if (track.durationSec < 30) {
            std::cout << "[Skipped - Too Short]: " << track.title << "\n";
            return false;
        }
        return TrackFilter::process(track);
    }
};

// --- ITERATOR + CHAIN OF RESPONSIBILITY IN ACTION ---
void runAudioPipeline(TrackIterator& it, TrackFilter& pipeline) {
    std::cout << "--- Processing Tracks Through Filter Pipeline ---\n";
    while (it.hasNext()) {
        const Track& track = it.next();
        // Iterator provides the item; Chain processes/filters the item
        if (pipeline.process(track)) {
            std::cout << "[Playing]: " << track.title << "\n";
        }
    }
}
\end{lstlisting}